\section{Fase 6, Servicio Systemd}

\textbf{Objetivo:} Confirmar que el grabador se inicia automáticamente con el sistema operativo y que se recupera de forma autónoma tras un fallo del proceso.

Si el servicio no está aún instalado, instálelo con los siguientes comandos antes de continuar:

\begin{lstlisting}[language=bash]
sudo cp /home/kraken/Signal-Recorder/signal-recorder.service \
    /etc/systemd/system/
sudo systemctl daemon-reload
sudo systemctl enable --now signal-recorder
\end{lstlisting}

% --- TEST 14 ---
\subsection{T-14, Servicio activo inmediatamente tras la instalación}

\begin{lstlisting}[language=bash]
systemctl status signal-recorder
journalctl -u signal-recorder -n 20
\end{lstlisting}

\begin{tcolorbox}[colback=gray!5, colframe=deyposcolor, title=Resultado esperado (T-14)]
La salida de \texttt{status} debe mostrar \texttt{\textbf{Active: active (running)}}. En las últimas líneas del journal deben ser visibles los mensajes de calibración del piso de ruido generados en el arranque.
\end{tcolorbox}

% --- TEST 15 ---
\subsection{T-15, El servicio sobrevive a un kill forzado}

\begin{lstlisting}[language=bash]
# Simula un crash matando el proceso principal
sudo kill -9 $(pgrep -f "python main.py")

# Espera 35 s (RestartSec=30 + margen de seguridad)
sleep 35

systemctl status signal-recorder
\end{lstlisting}

\begin{tcolorbox}[colback=gray!5, colframe=deyposcolor, title=Resultado esperado (T-15)]
Tras la espera, \texttt{status} vuelve a mostrar \texttt{\textbf{Active: active (running)}}. El journal refleja el ciclo completo: proceso terminado por señal 9 y, a continuación, una nueva secuencia de calibración del ruido de fondo.
\end{tcolorbox}

\begin{important}
    \textbf{Si el servicio no se recupera:} Es posible que \texttt{systemd} haya alcanzado el límite de reinicios (\texttt{StartLimitIntervalSec}). Ejecute: \texttt{sudo systemctl reset-failed signal-recorder \&\& sudo systemctl start signal-recorder}
\end{important}

% --- TEST 16 ---
\subsection{T-16, El servicio sobrevive un reinicio del sistema}

\begin{lstlisting}[language=bash]
sudo reboot
\end{lstlisting}

Tras el reinicio, conéctese de nuevo por SSH y compruebe inmediatamente el estado del servicio:

\begin{lstlisting}[language=bash]
systemctl status signal-recorder
\end{lstlisting}

\begin{tcolorbox}[colback=gray!5, colframe=deyposcolor, title=Resultado esperado (T-16)]
\texttt{\textbf{Active: active (running)}} sin ninguna intervención manual. El servicio ha arrancado automáticamente con el sistema gracias a \texttt{systemctl enable}.
\end{tcolorbox}

% --- TEST 17 ---
\subsection{T-17, El heartbeat registra el estado del servicio}

Tras al menos una hora de servicio activo, verifique el archivo de estado periódico:

\begin{lstlisting}[language=bash]
tail /home/kraken/Signal-Recorder/status.log
\end{lstlisting}

\begin{tcolorbox}[colback=gray!5, colframe=deyposcolor, title=Resultado esperado (T-17)]
Líneas con el siguiente formato:
\vspace{0.2cm}
\begin{center}
    \texttt{[2026-05-09 15:00:00] Red: OK | Servicio: \textbf{active} | Archivos pendientes: 4}
\end{center}
\vspace{0.2cm}
El campo \texttt{Servicio} debe mostrar \texttt{\textbf{active}}. Los valores \texttt{failed} o \texttt{unknown} constituyen un fallo de esta prueba.
\end{tcolorbox}
